%Please use LuaLaTeX or XeLaTeX
\documentclass[11pt,aspectratio=169]{beamer}

\usepackage[ngerman]{babel}
\usepackage{ragged2e}

\definecolor{ethblue}{RGB}{32, 91, 174}
\definecolor{mygrey}{RGB}{175, 175, 175}

\hypersetup{
  colorlinks=true,
  urlcolor=ethblue,
  linkcolor=ethblue,
  citecolor=ethblue
}

\title{Theoretische Informatik}
% \date[ABC 2021]{24. Sep 2026}
\date[]{29. Okt 2026}
\author{Olivér Kovács}
% \institute{Organizational Unit\\can spread over 2 lines}

\usetheme{eth}

\colorlet{titlefgcolor}{ETHBlue}
\colorlet{accentcolor}{ETHRed}

\begin{document}

%\def\titlefigure{elements/title-page-image}		% Default image
%\def\titlefigure{elements/title-page-image-43}	% Use this for 4:3 presentations

\titleframe

% \tocframe

% \section{Layout}

\begin{frame}{Midterm}
    Altklausuren: \href{https://exams.vis.ethz.ch/}{https://exams.vis.ethz.ch/}

    \vspace{1em}

    Stoff hat sich geändert:
    \begin{itemize}
    \item
        Kolmogorovkomplexität, Satz 3.1 nicht relevant.
    \item
        Relevant bis einschliesslich
        \begin{itemize}
            \item
                Vorlesung am 03. Novemeber,
            \item
                Blatt 6,
            \item
                Quiz 7.
        \end{itemize}
    \item Manches davon war früher nicht auf den Midterms!
    \end{itemize}
\end{frame}

\begin{frame}{Bemerkungen zur Serie}
    \begin{itemize}
    \item
        9 (b). Mathematische Angabe erheblich einfacher als Beschreibung des Konstruktionsverfahrens.
        \vspace{1em}
    \item
        9 (c). Falls konkreter Automat verlangt ist, einfach Bild zeichnen. :)
        \vspace{1em}
    \item
        10 (a). Für allgemeine Konstruktionen reicht ein Bild nicht.
        \vspace{1em}
    \item
        10 (a). Regularität von \(L^R\) muss bewiesen werden.
    \end{itemize}
\end{frame}

\begin{frame}{Quiz}
    Gibt es Fragen zum Quiz?
\end{frame}

\begin{frame}{Reduktionen}

    \justifying

    \textbf{Definition 5.3.}
    Die Sprache \(L_1\) ist auf \(L_2\) \textit{rekursiv reduzierbar},
    \(L_1 \le_{\text{R}} L_2\), falls
    \[
        L_2 \in \mathcal L_{\text{R}} \implies L_1 \in \mathcal L_{\text{R}}. 
    \]

    \vspace{1em}

    \textbf{Definition 5.4.}
    Die Sprache \(L_1\) ist auf \(L_2\) \textit{EE-reduzierbar},
    \(L_1 \le_{\text{EE}} L_2\), falls eine TM \(M\) existiert,
    welche die Funktion \(f_M\) berechnet, so dass gilt
    \[
        x \in L_1 \iff f_M(x) \in L_2.
    \]

    \textbf{Lemma 5.4.}
    Für jede Sprache \(L\) gilt \(L \le_{\text{R}} L^C\).

    \vspace{1em}

    Beweis?
\end{frame}

\begin{frame}{Diagonalsprache}

    \justifying

    \textbf{Definition.}
    Seien \(w_1, w_2, w_3, \ldots\) die Wörter über \(\{0, 1\}\) in kanonischer Ordnung
    und \(M_1, M_2, M_3, \ldots\) die Folge aller Turingmaschinen. Wir definieren die
    \textit{Diagonalsprache} \(L_{\text{diag}}\) als

    \[
        L_{\text{diag}} =
        \{ w \in \{0, 1\}^* \mid w = w_i \text{ und } M_i \text{ akzeptiert } w_i \text{ nicht} \}.
    \]

    \vspace{1em}

    \textbf{Satz 5.5.} \(L_{\text{diag}} \not\in \mathcal L_{\text{RE}}\).

    \vspace{1em}

    \textbf{Lemma 5.5.} \((L_{\text{diag}})^C \in \mathcal L_{\text{RE}}\).

    \vspace{1em}

    Beweis?
\end{frame}

\begin{frame}{Die universelle Sprache}

    \justifying

    \textbf{Definition 5.6.}
    Die \textit{universelle Sprache} ist die Sprache

    \[
        L_U = \{\text{Kod}(M)\#w \mid w \in \{0, 1\}^* \text{ und } M \text{ akzeptiert } w\}.
    \]

    \vspace{1em}

    \textbf{Satz 5.6.} \(L_{\text{U}} \in \mathcal L_{\text{RE}}\).

    \vspace{1em}

    \textbf{Satz 5.7.} \(L_{\text{U}} \not\in \mathcal L_{\text{R}}\).

    \vspace{1em}

    Beweis?
\end{frame}

\begin{frame}{Das Halteproblem}

    \justifying

    \textbf{Definition 5.6.}
    Das \textit{Halteproblem} ist das Entscheidungsproblem \((\{0, 1, \#\}, L_H)\), wobei
    \[
        L_H = \{\text{Kod}(M)\#w \mid w \in \{0, 1\}^* \text{ und } M \text{ hält auf } w\}.
    \]

    \vspace{1em}

    \textbf{Satz 5.8.} \(L_{\text{H}} \not\in \mathcal L_{\text{R}}\).

    \vspace{1em}

    Beweis?
\end{frame}

\begin{frame}{\(L_{\text{empty}}\)}

    \justifying

    \textbf{Definition.} Sei \(L_{\text{empty}}\) die Sprache
    
    \[
        L_{\text{empty}} = \{\text{Kod}(M) \mid L(M) = \varnothing\}.
    \]

    \vspace{1em}

    \textbf{Lemma 5.6.} \((L_{\text{empty}})^C \in \mathcal L_{\text{RE}}\).

    \vspace{1em}

    \textbf{Lemma 5.7.} \((L_{\text{empty}})^C \not\in \mathcal L_{\text{R}}\).

    \vspace{1em}

    Beweis?
\end{frame}

\begin{frame}{\(L_{\text{empty}}\)}

    \justifying

    \textbf{Korollar 5.4.} \(L_{\text{empty}} \not\in \mathcal L_{\text{R}}\).

    \vspace{1em}

    \textbf{Korollar 5.5.} Sei \(L_{\text{EQ}}\) die Sprache

    \[
        L_{\text{EQ}} = \{\text{Kod}(M)\#\text{Kod}(M') \mid L(M) = L(M')\}.
    \]

    Es gilt \(L_{\text{EQ}} \not\in \mathcal L_{R}\).

    \vspace{1em}

    Beweis?
\end{frame}

\setbeamertemplate{footline}{}

\begin{frame}{Aufgaben}
    \begin{enumerate}
        \justifying
        \item
            \textbf{Aufgabe 1.}
            Beschreibe informell eine NTM, welche die Sprache
            \[
                L_{\text{composite}} = \{ w = 0^n \mid n \in \mathbb N \text{ ist nicht prim}\}
            \]
            akzeptiert.
        \item
            \textbf{Aufgabe 2.}
            Wir betrachten die Sprache
            \[
                L_{\lambda + 1} = \{\text{Kod}(M) \mid \text{TM } M \text{ akzeptiert } \lambda \text{ und mindestens ein weiteres Wort}\}.
            \]
            Zeige, dass \(L_{\lambda + 1} \in \mathcal L_{\text{RE}}\).
        \item
            \textbf{Aufgabe 3.}
            Beweise, dass für jede Sprache \(L\) gilt
            \[
                L \in \mathcal L_{\text{RE}} \setminus \mathcal L_{\text{R}} \implies L^C \not\in \mathcal L_{\text{RE}}.
            \]
        \item
            \textbf{Aufgabe 4.}
            Gibt es eine Sprache \(L\), für die gilt
            \(L \not\in \mathcal L_{\text{RE}}\) und
            \(L^C \not\in \mathcal L_{\text{RE}}\)?

    \end{enumerate}
\end{frame}

\end{document}
